% Página 26 do PDF de Gabriel
\begin{frame}[plain,noframenumbering,label=original-26]
  \begin{tikzpicture}[remember picture,overlay,x=1bp,y=-1bp]
  \begin{scope}[shift={(current page.north west)}]
  \Fundo{26}
  \Texto{44.456}{41.894}{90.657}{8.518}{b}{34E0C8}{ANÁLISE DESCRITIVA}
  \Texto{43.039}{116.125}{63.147}{56.787}{b}{FFB014}{05}
  \Texto{43.039}{154.533}{155.397}{24.608}{b}{FFFFFF}{Distribuições}
  \Texto{43.039}{234.662}{46.285}{7.572}{b}{34E0C8}{EMASC 2026}
  \Texto{91.429}{234.662}{2.65}{7.572}{b}{FFB014}{\textbullet{}}
  \Texto{96.176}{234.662}{21.45}{7.572}{b}{34E0C8}{LAME}
  \Link{42.042}{135.586}{157.388}{25.418}{original-26}
  \end{scope}
  \end{tikzpicture}
\end{frame}


% Páginas 27 a 31 do PDF de Gabriel
\begin{frame}[plain,noframenumbering,label=original-27]
  \begin{tikzpicture}[remember picture,overlay,x=1bp,y=-1bp]
  \begin{scope}[shift={(current page.north west)}]
  \only<1>{\Fundo{27}}
  \only<2>{\Fundo{28}}
  \only<3>{\Fundo{29}}
  \only<4>{\Fundo{30}}
  \only<5>{\Fundo{31}}
  \Texto{40.175}{27.803}{198.275}{17.983}{b}{2C2084}{Construção do boxplot}
  \Texto{35.433}{63.506}{246.971}{8.518}{i}{737696}{Os quartis permitem um gráfico que resume a distribuição inteira.}
  \only<1>{\Texto{49.227}{215.166}{177.703}{8.518}{r}{000000}{Os mesmos sete valores, agora sobre um eixo.}}
  \Texto{37.21}{247.463}{40.499}{6.625}{b}{737696}{EMASC 2026}
  \Texto{80.653}{247.463}{2.319}{6.625}{r}{FFB014}{\textbullet{}}
  \Texto{85.913}{247.463}{117.821}{6.625}{r}{737696}{Liga Acadêmica Matemática \& Empresa}
  \Texto{206.682}{247.463}{2.319}{6.625}{r}{FFB014}{\textbullet{}}
  \Texto{211.943}{247.463}{13.25}{6.625}{r}{737696}{UFS}
  \Texto{397.549}{248.132}{7.367}{6.625}{b}{2C2084}{20}
  \Texto{406.02}{248.132}{10.315}{6.625}{r}{737696}{/ 40}
  \only<2>{\Texto{49.227}{215.166}{135.147}{8.518}{r}{000000}{Quartis: 25\% dos valores abaixo de}}
  \only<2>{\Texto{186.739}{215.166}{6.627}{8.518}{i}{000000}{Q}}
  \only<2>{\Texto{193.366}{216.162}{3.157}{5.679}{r}{000000}{1}}
  \only<2>{\Texto{197.022}{215.166}{161.936}{8.518}{r}{000000}{, 50\% abaixo da mediana e 75\% abaixo de}}
  \only<2>{\Texto{361.323}{215.166}{6.627}{8.518}{i}{000000}{Q}}
  \only<2>{\Texto{367.95}{216.162}{3.157}{5.679}{r}{000000}{3}}
  \only<2>{\Texto{371.606}{215.166}{2.368}{8.518}{r}{000000}{.}}
  \only<3>{\Texto{49.227}{215.166}{52.812}{8.518}{r}{000000}{A caixa vai de}}
  \only<3>{\Texto{104.406}{215.166}{6.627}{8.518}{i}{000000}{Q}}
  \only<3>{\Texto{111.033}{216.162}{3.157}{5.679}{r}{000000}{1}}
  \only<3>{\Texto{117.056}{215.166}{4.736}{8.518}{r}{000000}{a}}
  \only<3>{\Texto{124.16}{215.166}{6.627}{8.518}{i}{000000}{Q}}
  \only<3>{\Texto{130.787}{216.162}{3.157}{5.679}{r}{000000}{3}}
  \only<3>{\Texto{134.442}{215.166}{174.875}{8.518}{r}{000000}{; sua largura é o IQR e cobre os 50\% centrais.}}
  \only<4>{\Texto{49.227}{215.166}{182.728}{8.518}{r}{000000}{Os bigodes terminam no menor e no maior valor}}
  \only<4>{\Texto{234.321}{215.166}{26.329}{8.518}{b}{000000}{dentro}}
  \only<4>{\Texto{263.018}{215.166}{88.63}{8.518}{r}{000000}{das cercas: aqui, 2 e 8.}}
  \only<5>{\Texto{49.227}{215.327}{324.851}{8.518}{r}{000000}{Com 24 h no lugar de 8 h, o valor ultrapassa a cerca superior (8,5) e é marcado como}}
  \only<5>{\Texto{376.488}{215.327}{27.454}{8.518}{b}{000000}{atípico}}
  \only<5>{\Texto{403.941}{215.327}{2.368}{8.518}{r}{000000}{.}}
  \only<5>{\Texto{49.227}{226.286}{160.087}{8.518}{r}{000000}{A caixa e o IQR permanecem os mesmos.}}
  \end{scope}
  \end{tikzpicture}
\end{frame}


% Páginas 32 a 34 do PDF de Gabriel
\begin{frame}[plain,noframenumbering,label=original-32]
  \begin{tikzpicture}[remember picture,overlay,x=1bp,y=-1bp]
  \begin{scope}[shift={(current page.north west)}]
  \only<1>{\Fundo{32}}
  \only<2>{\Fundo{33}}
  \only<3>{\Fundo{34}}
  \Texto{40.175}{27.803}{230.787}{17.983}{b}{2C2084}{Construção do histograma}
  \Texto{35.433}{63.506}{304.297}{8.518}{i}{737696}{Para ver a forma da distribuição, contamos quantos valores caem em cada faixa.}
  \only<1>{\Texto{49.227}{209.496}{109.15}{8.518}{r}{000000}{Os sete valores sobre a reta.}}
  \only<1,2>{\Texto{37.21}{247.463}{40.499}{6.625}{b}{737696}{EMASC 2026}}
  \only<1,2>{\Texto{80.653}{247.463}{2.319}{6.625}{r}{FFB014}{\textbullet{}}}
  \only<1,2>{\Texto{85.913}{247.463}{117.821}{6.625}{r}{737696}{Liga Acadêmica Matemática \& Empresa}}
  \only<1,2>{\Texto{206.682}{247.463}{2.319}{6.625}{r}{FFB014}{\textbullet{}}}
  \only<1,2>{\Texto{211.943}{247.463}{13.25}{6.625}{r}{737696}{UFS}}
  \only<1,2>{\Texto{397.549}{248.132}{7.367}{6.625}{b}{2C2084}{21}}
  \only<1,2>{\Texto{406.02}{248.132}{10.315}{6.625}{r}{737696}{/ 40}}
  \only<2>{\Texto{49.227}{209.3}{357.083}{8.518}{r}{000000}{Faixas de mesma largura cobrem a amplitude dos dados; conta-se quantos valores caem em}}
  \only<2>{\Texto{49.227}{220.259}{40.938}{8.518}{r}{000000}{cada faixa.}}
  \only<3>{\Texto{49.227}{209.495}{357.083}{8.518}{r}{000000}{Uma barra por faixa, com altura igual à contagem. A largura da faixa é a amplitude dividida pelo}}
  \only<3>{\Texto{49.227}{220.454}{298.871}{8.518}{r}{000000}{número de faixas: faixas largas escondem o formato; estreitas produzem ruído.}}
  \only<3>{\Texto{37.21}{247.462}{40.499}{6.625}{b}{737696}{EMASC 2026}}
  \only<3>{\Texto{80.653}{247.462}{2.319}{6.625}{r}{FFB014}{\textbullet{}}}
  \only<3>{\Texto{85.913}{247.462}{117.821}{6.625}{r}{737696}{Liga Acadêmica Matemática \& Empresa}}
  \only<3>{\Texto{206.682}{247.462}{2.319}{6.625}{r}{FFB014}{\textbullet{}}}
  \only<3>{\Texto{211.943}{247.462}{13.25}{6.625}{r}{737696}{UFS}}
  \only<3>{\Texto{397.549}{248.131}{7.367}{6.625}{b}{2C2084}{21}}
  \only<3>{\Texto{406.02}{248.131}{10.315}{6.625}{r}{737696}{/ 40}}
  \end{scope}
  \end{tikzpicture}
\end{frame}


% Página 35 do PDF de Gabriel
\begin{frame}[plain,noframenumbering,label=original-35]
  \begin{tikzpicture}[remember picture,overlay,x=1bp,y=-1bp]
  \begin{scope}[shift={(current page.north west)}]
  \Fundo{35}
  \Texto{40.175}{27.803}{219.315}{17.983}{b}{2C2084}{Distribuições observadas}
  \Texto{35.433}{63.506}{280.285}{8.518}{i}{737696}{A forma da distribuição explica quando a média representa bem os dados.}
  \Texto{49.227}{203.827}{357.083}{8.518}{r}{000000}{Em distribuições assimétricas, a média é deslocada no sentido da cauda: na pontuação de}
  \Texto{49.227}{225.745}{225.616}{8.518}{r}{000000}{sono e de redes, média e mediana praticamente coincidem.}
  \Texto{37.21}{247.463}{40.499}{6.625}{b}{737696}{EMASC 2026}
  \Texto{80.653}{247.463}{2.319}{6.625}{r}{FFB014}{\textbullet{}}
  \Texto{85.913}{247.463}{117.821}{6.625}{r}{737696}{Liga Acadêmica Matemática \& Empresa}
  \Texto{206.682}{247.463}{2.319}{6.625}{r}{FFB014}{\textbullet{}}
  \Texto{211.943}{247.463}{13.25}{6.625}{r}{737696}{UFS}
  \Texto{397.549}{248.132}{7.367}{6.625}{b}{2C2084}{22}
  \Texto{406.02}{248.132}{10.315}{6.625}{r}{737696}{/ 40}
  \end{scope}
  \end{tikzpicture}
\end{frame}


% Página 36 do PDF de Gabriel
\begin{frame}[plain,noframenumbering,label=original-36]
  \begin{tikzpicture}[remember picture,overlay,x=1bp,y=-1bp]
  \begin{scope}[shift={(current page.north west)}]
  \Fundo{36}
  \Texto{40.175}{28.028}{215.844}{17.983}{b}{2C2084}{Comparação entre níveis}
  \Texto{35.433}{63.506}{225.276}{8.518}{i}{737696}{As mesmas medidas e gráficos permitem comparar grupos.}
  \Texto{49.227}{203.827}{357.083}{8.518}{r}{000000}{As médias de uso de IA diferem em centésimos de hora, enquanto o desvio padrão é de cerca}
  \Texto{49.227}{214.786}{86.296}{8.518}{r}{000000}{de 1,7 h: a diferença é}
  \Texto{138.092}{214.786}{117.881}{8.518}{b}{000000}{pequena diante da dispersão}
  \Texto{255.97}{214.786}{150.343}{8.518}{r}{000000}{. Além disso, os níveis correspondem a}
  \Texto{49.227}{225.745}{85.904}{8.518}{r}{000000}{faixas etárias distintas.}
  \Texto{37.21}{247.463}{40.499}{6.625}{b}{737696}{EMASC 2026}
  \Texto{80.653}{247.463}{2.319}{6.625}{r}{FFB014}{\textbullet{}}
  \Texto{85.913}{247.463}{117.821}{6.625}{r}{737696}{Liga Acadêmica Matemática \& Empresa}
  \Texto{206.682}{247.463}{2.319}{6.625}{r}{FFB014}{\textbullet{}}
  \Texto{211.943}{247.463}{13.25}{6.625}{r}{737696}{UFS}
  \Texto{397.549}{248.132}{7.367}{6.625}{b}{2C2084}{23}
  \Texto{406.02}{248.132}{10.315}{6.625}{r}{737696}{/ 40}
  \end{scope}
  \end{tikzpicture}
\end{frame}
